\scanpage{38}
% letterhead
\begin{center}
INSTITUT\\
DES HAUTES ÉTUDES SCIENTIFIQUES\\
\rule{2.3em}{0.4pt}
\end{center}

SGA 4 $\underline{\overline{\mathrm{XVIII}}}$ \inserted{(Deligne)}

Références pour SGA 5 $\underline{\overline{\mathrm{VII}}}$
\diagram{\(\left\{\begin{array}{l}\text{\unsure{Dèm|Dém} de 3.\deleted{\unsure{6}}5.3.}\\
\text{1.1\qquad b)}
\end{array}\right.\)}

\noindent\rule{\linewidth}{0.4pt}

$\underline{\overline{\mathrm{V}}}$\qquad dire quelques mots sur relations avec $\underleftarrow{\mathrm{lim}}$ et \({\underleftarrow{\mathrm{lim}}}^{(p)}\) :

a) Sur un topos de \unsure{la} forme $\widehat C$, les foncteurs $H^n$ sont les $\mathop{\underleftarrow{\mathrm{lim}}}\nolimits_C^{(n)}$

b) Il \unsure{n'y a} plus une telle interprétation pour un topos général, puisqu'on est obligé de prendre
\[
 H^{n}(\deleted{\unsure{E|C}}\widetilde C,F)\simeq
 \mathop{\underrightarrow{\mathrm{lim}}}_{\unsure{\mathcal U|U}\in\unsure{\mathcal R|R}}
 \mathop{\underleftarrow{\mathrm{lim}}}_{\mathcal U}
 \deleted{\unsure{H^n|F^n}}\unsure{F^{n}/\mathcal U_{\varphi}|F/\mathcal U_{\varphi}|F/C_{/\varphi}}
\]
\margin{si ça marche \ldots}

$U$ crible variable

On a un recouvrement de l'objet final de $\widehat C$ par les objets de $C$

\margin{$X\timesY$}
% archive sheet number: 37
